\subsection{How small a violation can be certified, and from which data}
\label{sec:limits}

Two questions bound what any procedure can do. How small a tolerance is
achievable, and when can the observed region of state space stand in for the
target domain?

\paragraph{A resolution limit.} Exact symmetry cannot be certified from noisy
data: a system with zero defect and one with a defect just above zero generate
almost the same data. The useful version of this statement is quantitative.

\begin{proposition}[Le~Cam bound on the certifiable tolerance]
\label{prop:lecam}
Consider model~\eqref{eq:model} with a fixed design whose empirical measure is
$\mu_N$, and a fixed generator $\xi = \xi_\theta \in \mathcal V$; write
$L_\xi$, $M^{(1)}_\xi$, $M^{(2)}_\xi$ and $M_\xi$ for the matrices $L_\theta$,
$M^{(1)}_\theta$, $M^{(2)}_\theta$ and $M_\theta$ of \Cref{lem:bilinear}.
Let $\varphi \in \{0,1\}$ be any rule
for certifying $\xi$ at tolerance $\delta$ whose false-certification rate is at
most $\alpha$, i.e.\ $\Prob_\beta(\varphi = 1) \le \alpha$ whenever
$\defect_\nu(\xi;F_\beta) > \delta$. Let $s_{\max}(\xi) = \max_{\Delta \neq 0}
\|L_\xi \Delta\|/\|F_\Delta\|_{\mu_N}$, attained at $\Delta^\star$ normalised
so that $\|F_{\Delta^\star}\|_{\mu_N} = 1$. Then for any $\beta^0$ with
$\defect_\nu(\xi;F_{\beta^0}) = 0$ and any $\delta < s_{\max}(\xi)/\|M_\xi
\Delta^\star\|$,
\begin{equation}
  \Prob_{\beta^0}(\varphi = 1)
  \;\le\; \alpha \;+\; \frac{\sqrt{N}}{2\sigma}\,
  \frac{\delta\,\|M_\xi \beta^0\|}{s_{\max}(\xi) - \delta\,\|M_\xi\Delta^\star\|} .
  \label{eq:lecam}
\end{equation}
Equivalently, attaining power $\gamma$ at $\beta^0$ requires
$\delta \ge 2\sigma(\gamma-\alpha)\,s_{\max}(\xi)\,/\,
\bigl(\sqrt{N}\,\|M_\xi\beta^0\| + 2\sigma(\gamma-\alpha)\,\|M_\xi\Delta^\star\|\bigr)$.
\end{proposition}

The bound is non-asymptotic and its constants are computable for the design at
hand, so it serves as a numerical benchmark for a single, pre-specified
generator in a fixed model class. It is a lower bound obtained through
Pinsker's inequality and one two-point alternative, and we do not claim it is
attained. \Cref{sec:resolution} compares it with the certificate: the two share
the $\sigma/\sqrt N$ scaling over two decades of $N$ at a fixed model class,
with a gap of a factor of about $8$ that increases with the size of the model
class. That comparison supports the scaling in $N$ and $\sigma$, not
optimality: neither the constant nor the dependence on the number of
coefficients $Q$ is addressed by the bound.

\paragraph{The price of incomplete coverage.} Suppose the data occupy a region
with distribution $\mu$ while the claim is about $\nu$. Both terms
$DF_\beta\,\xi_\theta$ and $D\xi_\theta\,F_\beta$ of every bracket, and hence
the bracket itself, lie in a fixed finite-dimensional space $\mathcal R$ of
fields (for polynomial classes, the polynomial fields up to the degree of a
bracket), so the worst-case discrepancy between the two measures is a pair of
generalised eigenvalues.

\begin{proposition}[Extrapolation factor]
\label{prop:kappa}
Let $G_\mu$ and $G_\nu$ be the Gram matrices of a basis of $\mathcal R$ under
$\mu$ and $\nu$, with $G_\nu$ positive definite, and set
$\kappa = \lambda_{\max}(G_\mu^{-1/2} G_\nu G_\mu^{-1/2})$ and
$\kappa' = \lambda_{\max}(G_\nu^{-1/2} G_\mu G_\nu^{-1/2})$, with $\kappa =
\infty$ when $G_\mu$ is singular. Then for every $q \in \mathcal R$,
$\|q\|^2_{\nu} \le \kappa\,\|q\|^2_{\mu}$ and
$\|q\|^2_{\mu} \le \kappa'\,\|q\|^2_{\nu}$, with equality attained, and
consequently, for every field $F$ in the class,
\begin{equation}
  \defect_\nu(\xi;F) \;\le\; \sqrt{\kappa\kappa'}\;\defect_\mu(\xi;F)
  \qquad\text{for every } \xi \in \mathcal V ,
\end{equation}
where $\defect_\mu$ is \Cref{def:defect} with $\nu$ replaced by $\mu$.
Moreover $\kappa = \infty$ if and only if some nonzero $q \in \mathcal R$
vanishes $\mu$-almost everywhere, i.e.\ the data lie in the zero set of a
nonzero field in $\mathcal R$, such as a possible bracket.
\end{proposition}

As a \emph{diagnostic}, $\kappa$ and $\kappa'$ are computable from the observed
states alone, before any symmetry analysis. In floating point $G_\mu$ is
treated as singular when its smallest eigenvalue is at most $10^{-10}$ times
its largest; no pseudoinverse or regularisation is used, and the factor
$\sqrt{\kappa\kappa'}$ is then reported as infinite. As an
\emph{explanation}, apply the proposition to the true field, assumed to lie
in the fitted class: generators whose defects measured on the data are at
most $\delta$ have true defects $\defect_\nu \le \sqrt{\kappa\kappa'}\,\delta$.
Let $r^\star$ be the dimension of the true symmetry algebra and
$\defect^\star$ the square root of the $(r^\star{+}1)$-th smallest
generalised eigenvalue of the true pencil $(C(\beta^\star), D(\beta^\star))$
on $\nu$ (for a system without symmetry, its smallest true defect). By the
Courant--Fischer theorem every subspace of dimension $r^\star + 1$ contains a
generator with true defect at least $\defect^\star$, so declaring more
generators than the true algebra has is possible only if
\begin{equation}
  \sqrt{\kappa\kappa'} \;\ge\; \defect^\star / \delta .
  \label{eq:kappacrit}
\end{equation}
This is a one-sided, population-level statement. Falling below the threshold
rules out declaring more generators than the true algebra has, and any
declaration for a system without symmetry. It does not rule out a declaration
of the right dimension tilted towards generators with true defect between
$\delta$ and $\sqrt{\kappa\kappa'}\,\delta$, nor errors from estimating the
field. Exceeding the threshold is a warning rather than a prediction, since
the bound is worst-case over $\mathcal R$. The infinite case is not exotic: a
single trajectory of the Hopf normal form converges onto its limit cycle, the
unit circle, and once the transient is discarded every field in $\mathcal R$
divisible by $x^2 + y^2 - 1$ vanishes on the data to within the eigenvalue
threshold, so the factor is reported as infinite. Single trajectories of the
Van der Pol and Lorenz systems give large but finite factors
(\STtrajectory).

\paragraph{Additive derivative noise and the plug-in defect.} It is natural
to expect that noise makes small singular values smaller. Under
model~\eqref{eq:model}, that is, with noise on a measured right-hand side, a
correctly specified class and exactly observed states, the opposite happens
near a true symmetry. Since $\E\|L_\theta \bhat\|^2 = \|L_\theta\beta^\star\|^2 +
\mathrm{tr}(L_\theta \Sigma L_\theta^{\!\top})$, where $\Sigma =
\mathrm{Cov}(\bhat)$, and likewise for the denominator, estimation error
inflates \emph{both} parts of equation~\eqref{eq:ratio}; for a direction close
to a true symmetry the numerator is near zero, so its relative inflation is
unbounded while the denominator's is not. The plug-in defect of a true
symmetry is therefore biased upward and all directions are pulled toward a
common noise-dominated value, so a fixed tolerance loses \emph{detections} as
the noise grows. For a direction whose true defect is well above zero the
argument gives no sign for the bias, since the common value may lie above or
below it, and a second effect pulls downward: the generator chosen to
\emph{minimise} the plug-in defect is selected for a small value (the
winner's curse), which can push the minimal plug-in defect below its true
value. On a system without symmetry, whether noise produces false
declarations therefore depends on how far the smallest true defect lies
above $\delta$; in our tests a threshold rule at $\delta$ declared a false
symmetry in at most $3.7\%$ of replications, and only when that defect was
within a factor of about two of $\delta$. Both effects are measured in
\Cref{sec:modelclass}, and the winner's curse at small samples and high noise
in \Cref{sec:tightness}. The argument concerns individual directions: on a
symmetric system the noise also tilts the estimated subspace away from the
true algebra, and a fixed tolerance can then declare generators whose true
defect exceeds it, which in our tests was the more frequent failure
(\Cref{sec:calibration}). \Cref{prop:simultaneous} removes the need to reason
about either the winner's curse or the tilt. The bias argument above holds
only under model~\eqref{eq:model}, with noise on the derivative and exact
states: noise in the \emph{states} acts quite differently, by attenuating
nonlinear terms, and can manufacture symmetry or hide it, depending on the
system (\Cref{sec:robustness}).
